\section{Билет 13. Затухающие колебания. Вывод уравнения затухающих колебаний.}

\begin{definition}
\textbf{Затухающими колебаниями} называются колебания, амплитуда которых уменьшается со временем из-за потери энергии на преодоление сил сопротивления среды или внутреннего трения. Этот процесс называется диссипацией энергии.
\end{definition}

\begin{theorem}[Вывод уравнения затухающих колебаний]
При движении тела в вязкой среде (жидкость или газ) на него действует сила сопротивления, которая при малых скоростях пропорциональна скорости:
\begin{equation}
    F_{\text{сопр}} = -r v = -r \dot{x}, \label{eq:friction_13}
\end{equation}
где $r$ --- коэффициент сопротивления (трения), зависящий от формы тела и вязкости среды. Знак «минус» указывает на противоположное направление силы относительно скорости.

Возвращающая (квазиупругая) сила имеет вид $F_{\text{упр}} = -kx$. Запишем второй закон Ньютона для одномерного движения:
\begin{equation}
    m \ddot{x} = F_{\text{упр}} + F_{\text{сопр}} = -kx - r\dot{x}.
\end{equation}
Перенесём все слагаемые в левую часть:
\begin{equation}
    m \ddot{x} + r \dot{x} + kx = 0.
\end{equation}
Разделим уравнение на массу $m$:
\begin{equation}
    \ddot{x} + \frac{r}{m} \dot{x} + \frac{k}{m} x = 0.
\end{equation}
Введём обозначения:
\begin{itemize}
    \item $\omega_0^2 = \frac{k}{m}$ --- квадрат собственной частоты колебаний (при отсутствии затухания);
    \item $2\beta = \frac{r}{m} \Rightarrow \beta = \frac{r}{2m}$ --- коэффициент затухания.
\end{itemize}
Подставляя их, получаем каноническое \textbf{уравнение затухающих колебаний}:
\begin{equation}
    \ddot{x} + 2\beta \dot{x} + \omega_0^2 x = 0. \label{eq:damped_ode}
\end{equation}
\end{theorem}

\begin{remark}
Уравнение \eqref{eq:damped_ode} является линейным однородным дифференциальным уравнением второго порядка с постоянными коэффициентами. Его характер зависит от соотношения между коэффициентом затухания $\beta$ и собственной частотой $\omega_0$. В курсе общей физики преимущественно рассматривается случай малого затухания $\beta < \omega_0$, когда система совершает колебания с постепенно убывающей амплитудой.
\end{remark}
